% ===================================================================
\section{Second and third laws: entropy, free energy}\label{sec:secondlaw}
% ===================================================================
\begin{multicols}{2}
\subsection{Statements and reversibility}\label{sec:2ndstatements}
\CORE{} \textbf{Kelvin--Planck:} no cyclic device can take heat from a single reservoir and convert it entirely to work. \textbf{Clausius:} no cyclic device can move heat from cold to hot with no other effect (no work input). A violation of either implies a violation of the other.

\textbf{Irreversibilities} (each generates entropy): heat transfer across a finite $\Delta T$, friction, unrestrained expansion, mixing of different substances, electrical resistance, inelastic deformation, chemical reaction away from equilibrium. A process is reversible only if both system \emph{and} surroundings can be restored with no net change anywhere.

\subsection{Entropy: definition and balance}\label{sec:entropy}
\CORE{}
\begin{formulabox}
\vspace{-6pt}
\begin{align}
\dd S &= \frac{\delta Q_{\text{rev}}}{T} \label{F:entropy:def}\\
\oint\frac{\delta Q}{T_b} &\le 0\quad\text{(Clausius inequality)} \label{F:entropy:clausius}\\
\Delta S_{\text{sys}} &= \sum\int\frac{\delta Q}{T_b}+\Sgen,\qquad \Sgen\ge0 \label{F:entropy:balance}
\end{align}
\end{formulabox}
$T_b$ = absolute temperature \emph{at the boundary where the heat crosses}. $S$ is a state function: to compute $\Delta S$ for an irreversible process, invent \emph{any} reversible path between the same end states and integrate $\delta Q_{\text{rev}}/T$ along it. The actual $Q/T$ of the irreversible path is smaller; the difference is $\Sgen$.

\textbf{Universe accounting:} $\Delta S_{\text{univ}}=\Delta S_{\text{sys}}+\Delta S_{\text{surr}}=\Sgen^{\text{total}}\ge0$ (= 0 only if reversible). Over a full cycle $\Delta S_{\text{device}}=0$, but $\Sgen=-\sum_k Q_k/T_k\ge0$ (heats into the device signed $+$) shows up in the reservoirs: \textbf{returning to the initial state does not make generation zero.}

\begin{formulabox}Entropy-change collection (state functions):
\begin{align}
\Delta S &= n\bar C_v\ln\frac{T_2}{T_1}+nR\ln\frac{V_2}{V_1} = n\bar C_p\ln\frac{T_2}{T_1}-nR\ln\frac{P_2}{P_1} \label{F:entropy:idealgas}\\
\Delta S &= mc\ln\frac{T_2}{T_1}\ \ \Big(\text{general: }\int\frac{C(T)}{T}\dd T\Big) \label{F:entropy:incompressible}\\
\Delta S &= \frac{\Delta H_{\text{tr}}}{T_{\text{tr}}}\ \ \text{(phase change at equilibrium } T_{\text{tr}}) \label{F:entropy:phase}\\
\Delta S_{\text{res}} &= \frac{Q_{\text{into res}}}{T_{\text{res}}}\ \ \text{(reservoir: } T \text{ constant)} \label{F:entropy:reservoir}\\
\Delta S_{\text{mix}} &= -nR\sum_i y_i\ln y_i\ \ \text{(ideal gases, same } T,P) \label{F:entropy:mixing}
\end{align}
\end{formulabox}
Notes: \eqref{F:entropy:idealgas} assumes constant $\bar C$; \eqref{F:entropy:phase} only at the equilibrium transition temperature (e.g.\ melting ice at \SI{0}{\celsius}: $333.4/273.15=\SI{1.221}{J/(g.K)}$)---for supercooled water freezing at \SI{-10}{\celsius}, build a three-step reversible path (warm, freeze at \SI{0}{\celsius}, cool). Mixing two samples of the \emph{same} gas gives $\Delta S_{\text{mix}}=0$ (Gibbs paradox resolved by indistinguishability).

\begin{demobox}\textbf{Entropy audit: hot copper into a lake.} \SI{0.500}{kg} Cu ($c=\SI{384}{J/(kg.K)}$, $C=\SI{192}{J/K}$) at \SI{100}{\celsius} into a lake at \SI{15}{\celsius} (reservoir).
Block: $\Delta S=192\ln(288.15/373.15)=\SI{-49.6}{J/K}$. Heat to lake: $Q=192(85)=\SI{16320}{J}$, $\Delta S_{\text{lake}}=16320/288.15=\SI{+56.6}{J/K}$. $\Sgen=\SI{7.0}{J/K}>0$ \checkmark.
\emph{Check:} numerical $\int C\,\dd T/T$ matches; sign of $\Sgen$ independent of which body is hotter.\src{DER}\end{demobox}

\begin{demobox}\textbf{Supercooled water freezes at \SI{-10}{\celsius}} (per gram; take $c_w=4.22$, $c_{\text{ice}}=\SI{2.06}{J/(g.K)}$ constant, $L_f=\SI{333.4}{J/g}$ at \SI{0}{\celsius}). The real process is irreversible, so build a reversible path: warm liquid to \SI{0}{\celsius} ($+0.1574$), freeze at \SI{0}{\celsius} ($-333.4/273.15=-1.2206$), cool ice to \SI{-10}{\celsius} ($-0.0768$): $\Delta S_{\text{sys}}=\SI{-1.1400}{J/(g.K)}$.
Heat released at \SI{-10}{\celsius} (Kirchhoff): $L=333.4-(4.22-2.06)(10)=\SI{311.8}{J/g}$; surroundings at \SI{263.15}{K}: $+1.1849$. $\Sgen=\SI{+0.045}{J/(g.K)}>0$: the system's entropy \emph{falls}, yet the process is spontaneous.\src{DER}\end{demobox}

\begin{demobox}\textbf{Mixing hot and cold water.} \SI{1.00}{kg} at \SI{80}{\celsius} + \SI{1.00}{kg} at \SI{20}{\celsius}, insulated, $c=\SI{4.18}{kJ/(kg.K)}$: $T_f=\SI{50}{\celsius}$. $\Delta S_{\text{hot}}=4180\ln(323.15/353.15)=\SI{-371.1}{J/K}$, $\Delta S_{\text{cold}}=\SI{+407.3}{J/K}$, total $\SI{+36.2}{J/K}$ generated by heat transfer across a finite $\Delta T$, even though no heat left the container.\src{DER}\end{demobox}

\subsection{T--S diagrams}\label{sec:ts}
For a \emph{reversible} process, $Q=\int T\,\dd S$ = area under the path on $T$--$S$ axes; a reversible cycle's enclosed area is net heat = net work (clockwise = engine with $T$ up, $S$ right). Reversible adiabatic = vertical line (isentropic); isothermal = horizontal line. For an irreversible path, the area is \emph{not} the heat.

\subsection{Statistical view}\label{sec:boltzmann}
\begin{formulabox}
\begin{equation}
S = k\ln\Omega \label{F:entropy:boltzmann}
\end{equation}
\end{formulabox}
$\Omega$ = number of microstates consistent with the macrostate. Doubling an ideal gas's volume multiplies $\Omega$ by $2^N$: $\Delta S=Nk\ln2=nR\ln2$, matching \eqref{F:entropy:idealgas}. ``Disorder'' is a loose metaphor; entropy counts accessible microstates (how energy and matter can be distributed), and some ordered-looking states have higher entropy.\src{EK}

\subsection{Free energies and spontaneity}\label{sec:freeenergy}
\CORE{}
\begin{formulabox}
\vspace{-6pt}
\begin{align}
A &\equiv U-TS,\qquad G\equiv H-TS \label{F:entropy:AG}\\
\Delta G &= \Delta H-T\Delta S\quad(\text{constant } T) \label{F:entropy:dG}\\
\dd G &\le 0\ \ (\text{const.\ }T,P;\ \text{no non-}P\dd V\text{ work}) \label{F:entropy:criterion}
\end{align}
\end{formulabox}
At constant $T,V$ use $\dd A\le0$. $-\Delta G$ = maximum non-expansion (e.g.\ electrical) work at constant $T,P$; $-\Delta A$ = maximum total work at constant $T$. ``Spontaneous'' means thermodynamically allowed, not fast.

\begin{reftab}{T:spont}{Sign logic for $\Delta G=\Delta H-T\Delta S$ (crossover $T^\ast=\Delta H/\Delta S$, if $\Delta H,\Delta S$ roughly constant).}
\small
\begin{tabularx}{\linewidth}{@{}ccY@{}}
\toprule
$\Delta H$ & $\Delta S$ & Favorable when\\
\midrule
$-$ & $+$ & all $T$\\
$+$ & $-$ & never (reverse is favorable)\\
$-$ & $-$ & low $T$ ($T<T^\ast$)\\
$+$ & $+$ & high $T$ ($T>T^\ast$)\\
\bottomrule
\end{tabularx}\end{reftab}

\begin{formulabox}Fundamental relations (closed, simple system):
\begin{align}
\dd U&=T\dd S-P\dd V, & \dd H&=T\dd S+V\dd P, \notag\\
\dd A&=-S\dd T-P\dd V, & \dd G&=-S\dd T+V\dd P. \label{F:entropy:fundamental}
\end{align}
\end{formulabox}
\EXT{} Maxwell relations: $(\partial T/\partial V)_S=-(\partial P/\partial S)_V$; $(\partial T/\partial P)_S=(\partial V/\partial S)_P$; $(\partial S/\partial V)_T=(\partial P/\partial T)_V$; $(\partial S/\partial P)_T=-(\partial V/\partial T)_P$.\src{DER}

\subsection{Chemical potential and equilibrium}\label{sec:chempot}
\begin{formulabox}
\vspace{-6pt}
\begin{align}
\mu_i &= \Big(\frac{\partial G}{\partial n_i}\Big)_{T,P,n_{j\ne i}},\quad \mu=\mu^\circ+RT\ln\frac{p}{p^\circ}\ \ (\text{ideal gas}) \label{F:entropy:mu}\\
\Delta_rG &= \Delta_rG^\circ+RT\ln Q \label{F:entropy:dGQ}\\
\Delta_rG^\circ &= -RT\ln K \label{F:entropy:K}
\end{align}
\end{formulabox}
Matter moves from high to low $\mu$; coexisting phases have equal $\mu$ for each species. $Q$ and $K$ are built from \textbf{dimensionless activities}: gas $p_i/p^\circ$ ($p^\circ=\SI{1}{bar}$), solute $c_i/c^\circ$ ($c^\circ=\SI{1}{mol/L}$), pure solid or liquid $=1$. $Q<K$: forward favorable; $Q=K$: equilibrium ($\Delta_rG=0$). Temperature dependence: \eqp{F:thermo:vanthoff}.

\subsection{Third law}\label{sec:thirdlaw}
\CORE{} As $T\to0$, the entropy of a \textbf{perfect crystal with a unique (nondegenerate) ground state} approaches zero (Planck's form of Nernst's theorem). Consequences: absolute entropies
\begin{equation}
S^\circ(T)=\int_0^T\frac{C_p}{T}\dd T+\sum\frac{\Delta H_{\text{tr}}}{T_{\text{tr}}} \label{F:entropy:absolute}
\end{equation}
are tabulated; $C_p\to0$ as $T\to0$; and \textbf{unattainability}: no finite sequence of processes reaches \SI{0}{K}.

\textbf{Residual entropy:} glasses and orientationally disordered crystals keep $S>0$ at $T\to0$---e.g.\ CO (up to $R\ln2$) and ice (Pauling estimate $R\ln\tfrac32=\SI{3.37}{J/(mol.K)}$).\src{EK}\src{DER}

\begin{trapbox}Elements in their reference state have $\Delta H_f^\circ=0$ but $S^\circ>0$ (graphite \num{5.7}, \ce{O2} \SI{205.2}{J/(mol.K)}).\src{CRC} Never set an element's entropy to zero in $\Delta_rS^\circ$ sums.\end{trapbox}

\EXT{} \textbf{Landauer limit:} erasing one bit dissipates at least $kT\ln2$ (\SI{2.87e-21}{J} at \SI{300}{K}). \textbf{Negative absolute temperature} occurs only in systems with bounded energy spectra (nuclear spins) with population inversion; such systems are ``hotter'' than any positive $T$. Not relevant to ordinary gases.
\end{multicols}
